% AAMAS 2027 author-side final writer pass (Round 3). Exploratory; not submission-ready.
% Frozen v15 evidence: 40/40 treatment grades, SHA-256 260db688e040671dc9f2c7a095573992dc6ceda3b271ebe1852a2e90d19b4586
\documentclass[sigconf,anonymous]{aamas}
\usepackage{booktabs}
\title{Selecting Earlier Task Records for Coding Agents: An Exploratory Study}
\author{Anonymous Author(s)}
\affiliation{\institution{Anonymous}}
\begin{document}
\begin{abstract}
We examine a source-only decision to retain historical coding-task records without access to the later target. We examine this distinction using historical issue records from SWE-ContextBench, fixed recency-based retrieval, and a custom binary selector that decides which records to retain and subsequently show to a worker. Among ten target tasks, a benchmark-annotated related record enters the retrieval set in eight, survives admission in two of those eight, and is exposed in one. On all ten matched target IDs, the separately run and test-path-filtered No-Memory condition resolves two tasks, while each tested memory condition resolves one; one No-Memory success depends on a sanitized-patch regrade. Most records contain truncated issue descriptions and reference-change paths rather than verified lessons from earlier agents. Independent worker runs also disagree despite identical prompts. The study traces removal of annotated related records, not the usefulness of those records or the causal effects of exclusion.
\end{abstract}
\keywords{coding agents, shared memory, historical task records, context selection, software engineering}
\maketitle

\section{Introduction}
An earlier Django issue reports that combining \texttt{QuerySet.only()} with \texttt{select\_related()} on a reverse \texttt{OneToOneField} includes fields that should have been deferred. A later issue describes a related query over a recursive \texttt{OneToOneField}: using these methods together raises a \texttt{FieldError}. The stored record names \texttt{django/db/models/sql/query.py}, but provides no tested repair procedure. Our stored record is a shortened issue description with a reference-change path, not a trace of how an earlier agent fixed the issue.

The later Django task has 114 earlier same-repository records. Fixed retrieval selects the three most recent; the related earlier issue is among them. A source-only selector retains it, and a task-aware selector exposes it to the worker. The worker does not resolve the later issue. This case distinguishes whether a record is available, whether it is delivered, and whether it is useful.

We evaluate a source-only admission decision that is not given the later task and therefore constrains what a later coding agent may see. A subsequent selection decision has more context but cannot recover a discarded record. The setting resembles information sharing across agents, although the evaluation uses independently reset coding-agent runs rather than a deployed team or a continuously updated memory service.

We compare No Memory, unrestricted sharing, random admission, custom binary admission, and that admission followed by a read filter. A diagnostic condition supplies a benchmark-annotated related record. The primary finding is a selection trace: early admission rejects six target-level occurrences of annotated related records, corresponding to only two distinct source decisions, and the read filter withholds almost all remaining records. None of the memory policies improves full task resolution in this pilot. Run-to-run variation limits causal interpretation.

\section{Related Work}
SWE-ContextBench studies reuse of context across related coding tasks, including previous trajectories, summaries, and retrieval methods \cite{zhu2026swecontext}. We use its Lite-derived task relationships but construct much simpler issue-based records, so its trajectory-based results are not replicated here. ReasoningBank derives reusable lessons from agent interactions \cite{ouyang2025reasoningbank}; our records have no observed source-agent outcome. MemGuard retains verifier signals to manage memory across its lifecycle \cite{wang2026memguard}; we do not evaluate verification, revision, or revocation.

Jev-Mem combines a System-One controller with structured memory construction and adaptive retrieval \cite{jiang2026jevmem}. We do \emph{not} run the published Jev-Mem system. We instead use the Jev System-One API to answer custom binary questions about storing and exposing issue records. These results must not be interpreted as an evaluation of Jev-Mem's graph memory or retrieval algorithm. ChainSWE evaluates linked repairs while retaining repository state \cite{jin2026chainswe}. We reset each worker to the prescribed repository state, separating explicit context delivery from changes accumulated in code.

Our narrower question is what happens when one source-only admission decision is reused across multiple later tasks. This reveals the repeated effect of one source-only admission choice across several later tasks. It is a descriptive account of where annotated related context becomes unavailable, not a new memory architecture, a useful-memory test, or a claim of superiority over earlier systems.

\section{Experimental Setting}
\subsection{Targets and records}
We use a frozen SWE-ContextBench Lite-derived release. Eligibility required at least one strictly earlier, different-pull-request source relation, at least five historical candidates, and at least two additional non-linked candidates with issue, path, or symbol overlap. Eleven eligible histories were selected by a deterministic SHA-256 ordering with seed 42. A Matplotlib target was excluded before the main treatments after a semantic no-op control passed its designated failing test. The remaining ten targets span Django, SymPy, and scikit-learn, with 13--114 earlier same-repository candidates per target. This is an eligibility-filtered, retrospectively assembled pilot rather than a representative or prospective sample of software tasks.

The eleven histories contain 211 distinct records; the union of the ten retained histories contains 188. Admission was applied once to all 211, including 23 records associated with the excluded Matplotlib history. Each record contains up to 320 characters of an earlier issue description, its issue-creation timestamp, and up to eight paths touched by the source reference patch. Direct inspection finds that \textbf{197 of 211} issue descriptions end in a truncation marker. All 211 have unknown source-agent outcomes, a generic precondition, and no verified repair lesson or source-agent trajectory. These are historical task records, not evidence of successful earlier agent behavior.

We require the source issue's released creation timestamp to precede the target timestamp and exclude the target's own reference change. This does not prove the source patch had been completed or observed by an agent at that time. The reference-change paths may have become available only after the source issue was opened, so the nominal temporal information boundary is imperfect.

Ten target IDs correspond to eight distinct resolving pull requests. \texttt{sympy-16342} and \texttt{sympy-16953} share SymPy PR 17526; \texttt{sympy-21309} and \texttt{sympy-9384} share PR 21309. The ten targets therefore do not represent ten independent repairs (Table~\ref{tab:pr-groups}).

\begin{table}[t]
\caption{Ten target IDs correspond to eight resolving pull requests. PR numbers are repository-specific.}
\label{tab:pr-groups}
\centering
\small
\begin{tabular}{llr}
\toprule
Repository & Target IDs & Target PR\\
\midrule
Django & 34176 & 16321\\
Django & 35356 & 18055\\
scikit-learn & 13771 & 13780\\
SymPy & 16342, 16953 & 17526\\
SymPy & 16946 & 17491\\
SymPy & 19235 & 19235\\
SymPy & 21203 & 21203\\
SymPy & 21309, 9384 & 21309\\
\bottomrule
\end{tabular}
\end{table}

\subsection{Selection policies}
\emph{Admission} is a decision to retain an earlier record in a shared store. We use a custom binary selector built with the TypeSafe SDK (version 0.7.2) and Jev System-One API model \texttt{jev-1.13.0}. It receives a source record without any later target and chooses SHARE or DO\_NOT\_SHARE. The SHARE criterion asks for useful, bounded, evidence-backed information with clear scope; the alternative covers weak support, specificity, staleness, redundancy, or possible negative transfer. This is a custom categorical prompt, not the Jev-Mem memory architecture. A decision is made once per source and reused wherever that source appears.

For each target, the fixed retrieval rule takes the three most recent earlier records \emph{before} applying admission decisions. Rejected records leave vacant slots; older records are not substituted. An optional \emph{exposure} selector then receives the current target issue, repository, and base commit together with each admitted retrieved record. It chooses EXPOSE or WITHHOLD using custom criteria concerning evidence match, relevance, redundancy, and possible negative transfer. All recorded decision rationales are empty; their specific reasons cannot be reconstructed from the output.

The conditions are A: No Memory; B: Share-All with fixed retrieval; C: random global admission; D: custom Jev-based admission; and E: that admission followed by custom Jev-based exposure. C retains the same number of records globally as D, not the same count or length for each target. F supplies the benchmark-annotated related record as a diagnostic; it is not a performance upper bound. A proposed deliberative LLM admission condition (G) was not run. The design omits Share-All followed by the read selector (Table~\ref{tab:factorial}), so it cannot isolate the independent read-filter effect under unrestricted admission or the admission-by-exposure interaction.

\begin{table}[t]
\caption{The unmeasured cell in the admission-by-exposure design.}
\label{tab:factorial}
\centering
\begin{tabular}{lcc}
\toprule
Admission & No read filter & Task-aware read\\
\midrule
Share-All & B & Not run\\
Custom Jev-based & D & E\\
\bottomrule
\end{tabular}
\end{table}

\begin{figure*}[t]
\centering
\begin{minipage}[t]{0.48\textwidth}
\fbox{\begin{minipage}{0.94\linewidth}\small
\textbf{Django: related record retained}\\[3pt]
Source \texttt{django-16910} (stored excerpt): \emph{On Django 4.2 calling only() with select\_related() ... All the fields from the related model are still included in the ...} [truncated]. Reference path: \texttt{django/db/models/sql/query.py}.
\begin{center}
source $\rightarrow$ top-three retrieval\\[-1pt]
$\downarrow$ admission: SHARE\\[-1pt]
$\downarrow$ exposure: EXPOSE\\[-1pt]
$\downarrow$ target \texttt{django-35356}\\[-1pt]
worker: unresolved (0/1 F2P)
\end{center}
\end{minipage}}
\end{minipage}\hfill
\begin{minipage}[t]{0.48\textwidth}
\fbox{\begin{minipage}{0.94\linewidth}\small
\textbf{SymPy: related record rejected}\\[3pt]
Source \texttt{sympy-16988} (stored excerpt): \emph{Intersection should remove duplicates ... The routine should give the same answer if duplicates are present ...} [truncated]. Reference path: \texttt{sympy/sets/sets.py}.
\begin{center}
source $\rightarrow$ top-three retrieval\\[-1pt]
$\downarrow$ admission: DO\_NOT\_SHARE\\[-1pt]
$\downarrow$ source excluded\\[-1pt]
two other records: WITHHOLD\\[-1pt]
target \texttt{sympy-16946}: E sees no card
\end{center}
\end{minipage}}
\end{minipage}
\caption{Two observed selection paths. The benchmark annotates the earlier issues as related to their targets, but this does not establish that either record improves a repair. The SymPy exclusion is one source decision reused across three later targets. Two other related records are excluded by top-three retrieval itself.}
\label{fig:paths}
\end{figure*}

\subsection{Worker and grading}
Each target starts from its prescribed repository state. B--E use the SWE-Edit baseline worker at scaffold revision \texttt{60ca673}, requested model alias \texttt{gpt-6-luna} with \texttt{xhigh} reasoning effort, at most 100 iterations, and a 1,800-second timeout. An immutable model checkpoint is not available in the recorded API metadata. The frozen SWE-ContextBench executable evaluator runs through an Apptainer adapter.

We report resolved tasks, the number of previously failing tests that pass (F2P), and the number of previously passing tests that remain passing (P2P). Partial F2P success is not full resolution. A and F were run earlier and then regraded after excluding test-file modifications; B--E were executed separately. This changed an important baseline result: the original A prediction patch for \texttt{sympy-21203} failed to apply and was unresolved, whereas its test-path-filtered sanitized patch applied and resolved the task (3/3 F2P, 124/124 P2P). Thus the reported A 2/10 includes a regraded success, not a success in the original raw evaluation. All 40 archived B--E submitted prediction patches contain no test-file changes, although some working trees do. Different execution and patch-normalization histories are a material comparability limitation.

\section{Results}
\subsection{Where related records are excluded}
The admission selector retains 123 of 211 source records and rejects 88. Among the 188 distinct records in the ten retained histories, it retains 111; random admission retains 108. The read selector receives 15 admitted-and-retrieved records, exposes one, and withholds 14.

The benchmark-annotated related record appears in the fixed three-record retrieval set for eight of ten targets. It survives admission for two of those eight and exposure for one. On \texttt{django-34176} and \texttt{sympy-19235}, the related source is absent from the top three, so neither selector can expose it. The six target-level admission exclusions are only two distinct decisions: source \texttt{sympy-16988} is rejected for three targets, and source \texttt{sympy-21055} for three others. These are not six independent policy mistakes. The saved decisions for these sources have confidence 0.28 and 0.21 respectively, but no rationales. The saved decisions contain no rationales, so we cannot determine why the selector rejected either SymPy record. The data show where exclusion occurred, not whether a related record would have helped.

\subsection{Executable outcomes}
All 40 planned B--E grades are present in the frozen archive. Every B--E policy is evaluated on the same ten targets, and A and F each have ten separately executed canonical grades. Table~\ref{tab:aggregate} reports the complete matched comparison.

\begin{table}[t]
\caption{Outcomes on all ten targets. F2P and P2P are aggregated test counts, not independent task trials. A and F were executed separately.}
\label{tab:aggregate}
\centering
\begin{tabular}{lrrr}
\toprule
Condition & Resolved & F2P & P2P \\
\midrule
A: No Memory & 2/10 & 9/48 & 2145/2145\\
B: Share-All & 1/10 & 7/48 & 2025/2145\\
C: Random admission & 1/10 & 8/48 & 2025/2145\\
D: Jev-based admission & 1/10 & 13/48 & 2025/2145\\
E: Admission+read & 1/10 & 7/48 & 2025/2145\\
F: Known-related & 1/10 & 8/48 & 2025/2145\\
\bottomrule
\end{tabular}
\end{table}

All four B--E policies resolve \texttt{sympy-19235}, which A also resolves. D and E expose no records on that target; E's only resolved target therefore receives no memory at all. The shared success is not evidence of useful memory. A additionally resolves \texttt{sympy-21203}; none of B--E does. D passes 13 of 48 F2P tests versus A's nine, largely because it passes four of six failing tests on \texttt{sympy-16946}, where A passes none. That target remains unresolved. D also passes two instead of one F2P tests on \texttt{django-34176}. These are partial test differences, not additional successful repairs.

\begin{table}[t]
\caption{F2P tests passed by target, with the number of designated failing tests in parentheses. On \texttt{sympy-21309}, B--E and F each lose 120 of 320 P2P tests, whereas A preserves all 320. Some target IDs share an underlying resolving pull request.}
\label{tab:targets}
\centering
\scriptsize
\setlength{\tabcolsep}{2pt}
\begin{tabular}{lrrrrrr}
\toprule
Target (F2P total) & A & B & C & D & E & F\\
\midrule
django-34176 (6) & 1 & 1 & 1 & 2 & 1 & 1\\
django-35356 (1) & 0 & 0 & 0 & 0 & 0 & 0\\
scikit-learn-13771 (5) & 2 & 2 & 2 & 2 & 2 & 2\\
sympy-16342 (11) & 1 & 1 & 1 & 1 & 1 & 2\\
sympy-16946 (6) & 0 & 0 & 0 & 4 & 0 & 0\\
sympy-16953 (11) & 0 & 0 & 0 & 0 & 0 & 0\\
sympy-19235 (1) & 1 & 1 & 1 & 1 & 1 & 1\\
sympy-21203 (3) & 3 & 1 & 2 & 2 & 1 & 1\\
sympy-21309 (2) & 1 & 1 & 1 & 1 & 1 & 1\\
sympy-9384 (2) & 0 & 0 & 0 & 0 & 0 & 0\\
\bottomrule
\end{tabular}
\end{table}

The known-related diagnostic F resolves one of ten tasks, versus two of ten for the sanitized/regraded A condition. Three A/F target pairs have different executable grades, including regressions. This is evidence of behavioral sensitivity, not evidence that related records improve repairs.

\begin{table}[t]
\caption{Number of historical records exposed to each target. Counts are not budget-matched across conditions.}
\label{tab:exposure}
\centering
\small
\begin{tabular}{lrrrr}
\toprule
Target & B & C & D & E\\
\midrule
django-34176 & 3 & 2 & 1 & 0\\
django-35356 & 3 & 2 & 2 & 1\\
scikit-learn-13771 & 3 & 1 & 3 & 0\\
sympy-16342 & 3 & 1 & 2 & 0\\
sympy-16946 & 3 & 1 & 2 & 0\\
sympy-16953 & 3 & 1 & 2 & 0\\
sympy-19235 & 3 & 3 & 0 & 0\\
sympy-21203 & 3 & 3 & 1 & 0\\
sympy-21309 & 3 & 3 & 1 & 0\\
sympy-9384 & 3 & 3 & 1 & 0\\
\midrule
Total & 30 & 20 & 15 & 1\\
\bottomrule
\end{tabular}
\end{table}

\subsection{Run variability and context quantity}
E exposes no records on nine of ten targets. On all nine, its saved task prompt is byte-for-byte identical to the corresponding A prompt. Nevertheless A resolves \texttt{sympy-21203} and E does not, and E loses 120 P2P tests on \texttt{sympy-21309} that A preserves. These are independent worker runs without controlled repeats. The discrepancies cannot be attributed to memory exposure, and differences among other policies are descriptive rather than causal effect estimates. The outcome-changing A regrade further prevents a causal interpretation of the one-task resolution gap.

Table~\ref{tab:exposure} shows that E exposes a record only on the unresolved Django target; its one resolved SymPy target receives none. Over ten targets, B exposes 30 records, C exposes 20, D exposes 15, and E exposes one. Their estimated memory-context token totals are 9,408, 6,267, 4,752, and 315, calculated from character counts divided by four rather than tokenizer measurements. Matching global admission counts does not match target-level exposure budgets. The evidence does not support a reliable monetary cost comparison.

\section{Discussion and Limitations}
Under fixed top-three retrieval, source-only admission removes six target-level occurrences of annotated related records, representing only two distinct source decisions. In this pilot the read filter reduces exposure to a single record. The same source decisions affect several later targets, illustrating a consequence of choosing what to store without knowing future tasks. This is an observation about access to records, not evidence that excluded records would have improved a repair.

Most issue descriptions are truncated, source-agent outcomes are unknown, and the records do not contain validated repair lessons. We cannot separate selection quality from the usefulness of the information selected. The published Jev-Mem system was not evaluated; custom criteria may have encouraged withholding, but empty rationales prevent attributing specific decisions to those criteria. These negative outcomes do not generalize to agent memory or memory governance as a whole.

The design lacks Share-All with task-aware read selection, so it cannot estimate independent read-filter effects under unrestricted admission. Top-three recency excludes two annotated related records before admission; random admission is not matched for exposed content or length. Ten target IDs cover eight underlying repairs and three repositories. A and F were run separately from B--E, and identical A/E prompts sometimes yield different grades. Controlled repeated runs and a common grading boundary are necessary before claiming policy effects.

A stronger study would first establish a positive control using a verified source-agent trajectory or independently checked lesson that reproducibly helps a later worker. It would then compare relevance-based retrieval and exposure-matched selection, add the missing admission-by-exposure condition, and repeat worker runs. The present study does not simulate memory revision, revocation, or continuing organizational learning.

\section{Conclusion}
We evaluated early retention and later exposure of historical issue records across ten coding-agent tasks. A related record reaches the worker in only one of eight cases where fixed retrieval makes it available. No tested memory policy improves full task resolution over the separately executed No-Memory baseline. The experiment traces removal of annotated related records but cannot establish their usefulness or whether exclusion caused the observed worker outcomes.

\bibliographystyle{ACM-Reference-Format}
\bibliography{references}
\end{document}